\section{符号说明}

本文统一用大写下标字母表示检测点，用向量记号表示平面点。为避免与最小包围圆半径混淆，干扰源的有效接收半径记作 $R$，而定位区域的\textbf{最小包围圆半径}一律记作 $\varrho(P)$（研究记录中曾写作 $R(P)$，本文正文不采用该写法）。

\begin{table}[H]
  \centering
  \begingroup
  \renewcommand{\arraystretch}{1.45}
  \begin{tabularx}{\textwidth}{C{0.16\textwidth}ZC{0.20\textwidth}}
    \toprule[1.2pt]
    {\heifont\bfseries 符号} & {\heifont\bfseries 定义} & {\heifont\bfseries 单位} \\
    \midrule[0.5pt]
    $\Omega$ & 目标源分布圆域，以原点为心、半径 1800 米的闭圆盘 & 米 \\
    $g$ & 某个干扰源的真实位置，运行时未知 & 米 \\
    $s_i$ & 第 $i$ 个检测点坐标 & 米 \\
    $\theta_i$ & 检测点 $s_i$ 处测得的示向度 & 度 \\
    $\varepsilon$ & 测向误差半宽，主模型取 $1^\circ$ & 度 \\
    $W_i$ & 第 $i$ 次观测对应的正向测角角域（两个半平面之交） & -- \\
    $P$ & 多次观测交会得到的定位区域，闭凸多面集 & -- \\
    $v_i$ & 定位区域 $P$ 的第 $i$ 个顶点 & 米 \\
    $D(P)$ & 定位区域直径，$\max_{x,y\in P}\lVert x-y\rVert$ & 米 \\
    $\varrho(P)$ & 定位区域的最小包围圆半径，$\min_c\max_{x\in P}\lVert x-c\rVert$ & 米 \\
    $\mathcal{C}(P)$ & 保证清除位置集合，$\bigcap_i\overline{B}(v_i,20)$ & -- \\
    $R$ & 单个干扰源的有效接收半径，取值于 $[1000,1500]$ & 米 \\
    $R_0$ & 有效接收半径下界，$R_0=1000$ & 米 \\
    $r$ & 首次检测点与真实源之间的距离，未知 & 米 \\
    $q$ & 第二个检测点或下一步动作位置，$q=s+au+bv$ & 米 \\
    $a,\ b$ & $q$ 相对首次示向度轴的纵向、横向位移 & 米 \\
    $B$ & 本步允许的移动预算 & 米 \\
    $h$ & 横向分离阈值，用于排除共线退化观测 & 米 \\
    $F(q)$ & 对候选点 $q$ 的最坏后验最小包围圆半径 & 米 \\
    $U(q)$ & $F(q)$ 的数值保守外包上界，$U(q)\ge F(q)$ & 米 \\
    $T$ & 整场任务的虚拟总时间 & 秒 \\
    $L$ & 全部动作之间的累计连续路程 & 米 \\
    $M$ & 整场检测次数 & 次 \\
    $S$ & 检测频道切换次数 & 次 \\
    $C,\ F$ & 清除成功次数、清除失败次数 & 次 \\
    $\tau$ & 源类型标识，取"全向"或"定向" & -- \\
    $n$ & 定向源的发射方向单位向量 & -- \\
    $Q$ & 覆盖证明中使用的凸子区域 & -- \\
    $S_Q$ & 用于排除凸子区域 $Q$ 的实际无信号站点集合 & -- \\
    \bottomrule[1.2pt]
  \end{tabularx}
  \endgroup
\end{table}

表中 $T$ 的账本形式在第三、四问中展开为
\begin{equation}
  T=\frac{L}{5}+5M+S+5C+3F ,
  \label{eq:time-ledger}
\end{equation}
其中 $L/5$ 为移动时间（速度为 5 米/秒），$5M$ 为检测时间（每次 5 秒），$S$ 为频道切换时间（每次 1 秒），$5C$ 与 $3F$ 分别为成功清除与失败清除的代价。该式来自题目附件给出的计时规则，不含机器狗自身的计算时间；程序运行时间另作硬约束处理。
